\chapterbackmatter{Solutions to Supplementary Exercises}

\begin{enumerate}
  \SupplementarySolution{1}{Cancellation Behind the Witten Index}
  Fix an eigenvalue \(E>0\), and restrict all operators to its
  eigenspace.  There
  \[
    q_E\equiv \frac{Q}{\sqrt E}
  \]
  is an involution: \(q_E^2=\mathbf1\).  Since
  \(\{q_E,(-1)^F\}=0\), it maps every bosonic eigenvector bijectively to
  a fermionic eigenvector of the same energy; its inverse is itself.
  Therefore
  \[
    \operatorname{Tr}_{E}(-1)^F=0
  \]
  for every \(E>0\).  The Boltzmann factor is constant on each such
  eigenspace, so the cancellation also holds in
  \(\operatorname{Tr}(-1)^F e^{-\beta H}\).  Only vectors with
  \(E=0\), on which \(Q\) may vanish without providing a partner, can
  contribute.  Thus the index is the number of bosonic supersymmetric
  ground states minus the number of fermionic ones and is independent of
  \(\beta\).

  \SupplementarySolution{2}{A Vacuum Escaping to Infinity}
  The auxiliary equation is controlled by
  \[
    f_\epsilon'(X)=f+\epsilon X.
  \]
  For every \(\epsilon\ne0\) there is one supersymmetric vacuum,
  \[
    X_{\rm vac}=-\frac f\epsilon,\qquad V_{\rm vac}=0.
  \]
  Its small fluctuations form a massive chiral multiplet of mass
  \(\lvert\epsilon\rvert\), and the vacuum contributes \(+1\) to the
  index.  As \(\epsilon\to0\), however,
  \(\lvert X_{\rm vac}\rvert\to\infty\).  At \(\epsilon=0\) the
  potential is the positive constant \(V=\lvert f\rvert^2\), so there is
  no zero-energy state and supersymmetry is broken.  Index invariance
  assumes that varying parameters does not alter the large-field
  asymptotics and does not let normalizable states enter from or escape
  to infinity.  That hypothesis fails here, exactly accounting for the
  change.

  \SupplementarySolution{3}{Normalized Goldstino Direction}
  At a stationary vacuum the tree-level fermion mass matrix annihilates
  the vector
  \[
    v_n=\sqrt2\,\mathcal F_{n0},\qquad
    v_A=iD_{A0}.
  \]
  The same statement holds nonperturbatively for the one-particle pole
  residues of the corresponding fermion fields.  Define
  \[
    \mathcal N^2
      =2\sum_n\lvert\mathcal F_{n0}\rvert^2
       +\sum_A D_{A0}^2 .
  \]
  Up to an overall Majorana phase, the normalized left-handed goldstino
  is therefore
  \[
    G_L=\frac{
      \sqrt2\sum_n\mathcal F_{n0}\psi_{nL}
      +i\sum_A D_{A0}\lambda_{AL}}
      {\mathcal N}.
  \]
  Section 29.2 gives
  \[
    \rho_{\rm VAC}
      =\sum_n\lvert\mathcal F_{n0}\rvert^2
       +\frac12\sum_A D_{A0}^2
      =\frac{\mathcal N^2}{2}.
  \]
  Thus the norm of the goldstino direction is fixed directly by the
  positive vacuum energy.

  \SupplementarySolution{4}{Holomorphic Spurion Selection Rule}
  Give \(\Phi\) charge \(+1\) under a bookkeeping \(U(1)\), and assign
  \(\mu_r\) charge \(-r\).  Independently choose \(R(\Phi)=r_\Phi\);
  since a superpotential has \(R=2\),
  \[
    R(\mu_r)=2-r\,r_\Phi.
  \]
  A Wilsonian correction must be holomorphic in \(\Phi,\mu_2,\mu_3\),
  have mass dimension three, and have both spurionic charges of the
  original action.  Analyticity at \(\mu_2=\mu_3=0\) excludes negative
  or fractional powers of the couplings.  The simultaneous charge and
  dimension equations then leave only the monomials
  \(\mu_2\Phi^2\) and \(\mu_3\Phi^3\).  Their coefficients cannot contain
  a nonconstant analytic neutral function of the couplings, because no
  dimensionless holomorphic neutral ratio is regular at the free point.
  Matching that point fixes the coefficients to their tree values.
  Perturbative effects may renormalize the Kahler potential but not this
  Wilsonian superpotential.

  \SupplementarySolution{5}{Supersymmetric One-Loop Gauge Coefficient}
  In a convention with
  \[
    \mu\frac{dg}{d\mu}=-\frac{b_0}{16\pi^2}g^3,
  \]
  gauge bosons contribute \(11C_1/3\), a Majorana gaugino contributes
  \(-2C_1/3\), a left-handed Weyl fermion contributes
  \(-2T(R)/3\), and a complex scalar contributes \(-T(R)/3\).
  The vector multiplet therefore contributes \(3C_1\), while each
  chiral multiplet contributes \(-T(R)\).  Summing all chiral
  representations gives
  \[
    b_0=3C_1-C_2.
  \]
  For \(SU(N_c)\), \(C_1=N_c\).  Each fundamental or antifundamental has
  \(T=1/2\), so \(N_f\) flavor pairs give \(C_2=N_f\), and
  \[
    b_0=3N_c-N_f.
  \]
  This is positive precisely in the asymptotically free range.

  \SupplementarySolution{6}{Instanton Weight as a Chiral Parameter}
  A Euclidean configuration of winding number \(\nu>0\) has
  \[
    I_E=\frac{8\pi^2\nu}{g^2}-i\nu\theta.
  \]
  Its semiclassical weight is consequently
  \[
    e^{-I_E}
      =\exp\left(-\frac{8\pi^2\nu}{g^2}+i\nu\theta\right)
      =\exp(2\pi i\nu\tau),
    \qquad
    \tau=\frac{4\pi i}{g^2}+\frac{\theta}{2\pi}.
  \]
  A Wilsonian \(\mathcal F\)-term is holomorphic in the chiral
  parameter \(\tau\), so positive-winding instantons supply positive
  powers of \(e^{2\pi i\tau}\).  Negative powers would grow like
  \(e^{+8\pi^2/g^2}\) as \(g\to0\), contradicting the weak-coupling
  expansion.  Anti-instantons instead depend on \(\tau^*\) and
  contribute to the conjugate \(\mathcal F\)-term, not to the same
  holomorphic function.

  \SupplementarySolution{7}{Mass-Deformed ADS Vacua}
  Put \(d=N_c-N_f\) and define
  \[
    A=\left(\frac{\Lambda^{3N_c-N_f}}{\det M}\right)^{1/d}.
  \]
  Differentiating \(dA+\operatorname{Tr}(mM)\) gives
  \[
    m=A\,M^{-1},\qquad M=A\,m^{-1}.
  \]
  Taking determinants and using the definition of \(A\) yields
  \[
    A^{N_c}=\det(m)\Lambda^{3N_c-N_f}.
  \]
  Hence there are \(N_c\) choices of phase and
  \[
    M=m^{-1}
      \left[\det(m)\Lambda^{3N_c-N_f}\right]^{1/N_c}.
  \]
  At a stationary point \(f_{\rm eff}=N_cA\).  Identifying
  \(A=\Lambda_{\rm low}^3\) gives the holomorphic threshold relation
  \[
    \Lambda_{\rm low}^{3N_c}
      =\det(m)\Lambda^{3N_c-N_f}.
  \]
  The \(N_c\) roots are the \(N_c\) supersymmetric vacua of the pure
  low-energy gauge theory.

  \SupplementarySolution{8}{Abelian Field-Strength \(\mathcal F\)-Term}
  In the chapter's four-component notation,
  \[
    W_L=\lambda_L
      +\frac12\gamma^\mu\gamma^\nu\theta_L f_{\mu\nu}
      -i\theta_LD
      +(\theta_L^{\mathsf T}\epsilon\theta_L)
       \sl{\partial}\lambda_R.
  \]
  Multiplying the two Grassmann-odd spinor superfields and reducing the
  gamma-matrix product gives
  \[
    -[W_L^{\mathsf T}\epsilon W_L]_{\mathcal F}
      =-2(\bar\lambda_R\sl{\partial}\lambda_R)
       -\frac12f_{\mu\nu}f^{\mu\nu}
       +\frac{i}{4}\epsilon^{\mu\nu\rho\sigma}
        f_{\mu\nu}f_{\rho\sigma}+D^2.
  \]
  Its real part contains the Maxwell, Majorana kinetic, and auxiliary
  terms.  Its imaginary part contains the theta density together with a
  derivative from the axial gaugino bilinear.  In particular,
  \[
    -\frac12\operatorname{Re}
      [W_L^{\mathsf T}\epsilon W_L]_{\mathcal F}
    =-\frac14f^2-\frac12\bar\lambda\sl{\partial}\lambda
      +\frac12D^2.
  \]

  \SupplementarySolution{9}{Metric from the \(N=2\) Holomorphic Function}
  Differentiating the Kahler potential gives
  \[
    K_{aa^*}
      =\frac{1}{4\pi}\operatorname{Im}h'(a)
      =\operatorname{Re}T(a).
  \]
  This is the common coefficient of the scalar, two fermion, and Abelian
  gauge kinetic terms required by \(N=2\) supersymmetry.  Local unitarity
  therefore requires
  \[
    \operatorname{Im}h'(a)>0.
  \]
  Perturbatively,
  \(h'(a)\sim(2i/\pi)\ln(a/\Lambda)\), up to an additive constant.
  Its imaginary part decreases without bound when \(\lvert a\rvert\) is
  taken sufficiently small.  Because a negative kinetic coefficient is
  impossible, the one-loop logarithm cannot describe the origin; new
  singularities and a dual weakly coupled description must intervene.

  \SupplementarySolution{10}{First-Order Abelian Dualization}
  Add a real multiplier \(\widetilde V\) so that its field strength
  \(\widetilde W_L\) couples linearly to the unconstrained \(W_L\).
  The \(W_L\)-dependent part of the action is quadratic, and stationarity
  gives
  \[
    W_L=-\frac{\widetilde W_L}{h'(\Phi)}.
  \]
  Substitution reverses the quadratic coefficient.  Defining
  \[
    \Phi_D=h(\Phi),\qquad h_D(\Phi_D)=-\Phi
  \]
  preserves the original \(N=2\) form because
  \[
    h_D'(\Phi_D)=-\frac{1}{h'(\Phi)}.
  \]
  Thus \(a_D=h(a)\).  Since
  \(T=h'/(4\pi i)\), the dual coefficient is
  \[
    T_D=\frac{h_D'}{4\pi i}
      =-\frac{1}{4\pi i\,h'}
      =\frac{1}{16\pi^2T}.
  \]
  Integrating out \(\widetilde V\) instead simply restores the original
  Bianchi identity, proving equivalence of the two formulations.

  \SupplementarySolution{11}{Generators of the Integral Duality Group}
  Electric--magnetic interchange and an integral theta shift act as
  \[
    S=\begin{pmatrix}0&1\\-1&0\end{pmatrix},
    \qquad
    T_b=\begin{pmatrix}1&0\\b&1\end{pmatrix},
    \quad b\in\mathbb Z.
  \]
  Both matrices have determinant one, so every product has integral
  entries and unit determinant; they generate \(SL(2,\mathbb Z)\).
  For
  \[
    M=\begin{pmatrix}n&m\\k&l\end{pmatrix},
    \qquad
    \binom{a'}{a_D'}=M\binom{a}{a_D},
  \]
  set \(\tau=h'(a)=da_D/da\).  Differentiation gives
  \[
    \tau'=\frac{da_D'}{da'}
      =\frac{k+l\tau}{n+m\tau}.
  \]
  The identity
  \(\operatorname{Im}\tau'
   =\operatorname{Im}\tau/\lvert n+m\tau\rvert^2\)
  shows that the positive kinetic region is preserved.

  \SupplementarySolution{12}{Theta Shift and the Dyon Lattice}
  The shift \(a_D\to a_D+ba\) is represented by \(T_b\) from the
  preceding solution.  To keep
  \(na+ma_D\) invariant, the integer row of charges must transform with
  the inverse matrix:
  \[
    (n,m)\longrightarrow(n,m)T_b^{-1}=(n-bm,m).
  \]
  Therefore the central charge and
  \(M_{\rm BPS}=\sqrt2\lvert na+ma_D\rvert\) are unchanged; only the
  labels assigned in the shifted electric frame differ.  Since the real
  part of \(a_D/a\) determines the electric charge induced on a magnetic
  state, increasing the theta angle by \(2\pi b\) changes a unit
  monopole's electric charge by \(be\).  Relabeling \(n\) by
  \(n-bm\) leaves the full dyon lattice invariant, which is the Witten
  effect and theta periodicity together.

  \SupplementarySolution{13}{Weak-Coupling Monodromy at Infinity}
  When \(u\) goes once counterclockwise around a sufficiently large
  circle,
  \[
    \sqrt{2u}\longrightarrow-\sqrt{2u},\qquad
    \ln(2u/\Lambda^2)\longrightarrow
      \ln(2u/\Lambda^2)+2\pi i.
  \]
  Hence \(a\to-a\).  Substituting both changes in the asymptotic formula
  for \(a_D\) gives
  \[
    a_D\longrightarrow 2a-a_D.
  \]
  In the ordering used in Chapter 29,
  \[
    \binom a{a_D}\longrightarrow
    M_\infty\binom a{a_D},
    \qquad
    M_\infty=
      \begin{pmatrix}-1&0\\2&-1\end{pmatrix}.
  \]
  Its determinant is one, as required for an allowed integral duality
  transformation.

  \SupplementarySolution{14}{Monodromy of a Massless Monopole}
  Encircling \(u_0\) once sends
  \[
    \ln(a_D/\Lambda_0)\longrightarrow
      \ln(a_D/\Lambda_0)+2\pi i,
  \]
  while the analytic coordinate \(a_D\) returns to itself.  The
  logarithmic term in \(a\) consequently changes by
  \[
    \frac{i}{\pi}a_D(2\pi i)=-2a_D.
  \]
  Therefore
  \[
    a\longrightarrow a-2a_D,\qquad a_D\longrightarrow a_D,
  \]
  or
  \[
    M_{\rm mono}=
      \begin{pmatrix}1&-2\\0&1\end{pmatrix}.
  \]
  The invariant vanishing combination is \(a_D\), confirming that the
  light state has electric--magnetic charge vector \((0,1)\).

  \SupplementarySolution{15}{Factorization of Seiberg--Witten Monodromies}
  The ordered product is
  \[
    M_+M_-=
    \begin{pmatrix}1&-2\\0&1\end{pmatrix}
    \begin{pmatrix}3&-2\\2&-1\end{pmatrix}
    =
    \begin{pmatrix}-1&0\\2&-1\end{pmatrix}
    =M_\infty.
  \]
  Thus a large contour can be deformed into contours around the two
  finite singularities with the stated ordering.  At \(u=+1\), the
  vanishing period is \(a_D\), so the massless state has charge
  \((n,m)=(0,1)\).  At \(u=-1\), it is \(a_D-a\), corresponding to
  \((n,m)=(-1,1)\).  Their charges have unit symplectic pairing, so no
  single electric frame makes both particles elementary everywhere.

  \SupplementarySolution{16}{Duality Covariance of the BPS Spectrum}
  Write the period column as \(v=(a,a_D)^{\mathsf T}\) and the charge row
  as \(q=(n,m)\).  Since the central charge is proportional to \(qv\),
  the transformation \(v\to Mv\) must be accompanied by
  \[
    q\longrightarrow q'=qM^{-1}.
  \]
  An \(SL(2,\mathbb Z)\) matrix has an integral inverse, so integer
  electric and magnetic charges remain integral.  Then
  \[
    q'v'=qM^{-1}Mv=qv,
  \]
  proving exact invariance of \(Z\) and
  \(M_{\rm BPS}=\lvert Z\rvert/2\).  The transformation changes the
  electric--magnetic description of a state, not its physical mass.

  \SupplementarySolution{17}{Wilsonian Versus 1PI Holomorphy}
  A Wilsonian action integrates out modes only above a nonzero sliding
  scale \(\lambda\).  Its couplings are therefore insensitive to
  singular propagation of massless particles at momenta below
  \(\lambda\), and its local \(\mathcal F\)-terms retain ordinary
  holomorphy in chiral fields and spurions.  The 1PI action integrates
  over all momenta.  With interacting massless fields, infrared poles
  can turn nonlocal expressions into terms that mimic nonholomorphic
  corrections to a local superpotential.  If the theory has a mass gap,
  or if all particles omitted from the effective description are
  massive and the selected gauge-invariant fields include every
  massless mode, the infrared ambiguity is absent.  Under those
  conditions the 1PI and Wilsonian effective superpotentials may be
  identified; otherwise the Wilsonian object is the reliable one for
  holomorphic selection rules.

  \SupplementarySolution{18}{Anomaly-Free \(R\) Charge in SQCD}
  Let \(R(Q)=R(\bar Q)=r\), while the gaugino has charge \(+1\).
  The mixed anomaly coefficient is
  \[
    N_c+N_f(r-1),
  \]
  because the \(2N_f\) matter Weyl fermions each have charge \(r-1\)
  and fundamental index \(1/2\).  Cancellation gives
  \[
    r=1-\frac{N_c}{N_f}.
  \]
  Consequently
  \[
    R(M)=2\left(1-\frac{N_c}{N_f}\right),\qquad
    R(B)=R(\bar B)=N_c\left(1-\frac{N_c}{N_f}\right).
  \]
  The fermionic component of any chiral composite has \(R\) charge one
  less than its scalar component.  These assignments are exact symmetry
  data even when the gauge coupling becomes strong.

  \SupplementarySolution{19}{Uniqueness of the ADS Superpotential}
  Flavor symmetry permits dependence only on \(\det M\), where
  \(M^i{}_j=\bar Q_jQ^i\).  Write
  \[
    f_{\rm np}=C\Lambda^p(\det M)^q.
  \]
  The anomaly-free \(R\) charge from the preceding solution requires
  \(2N_fq(1-N_c/N_f)=2\), so
  \(q=-1/(N_c-N_f)\).  Since \(\det M\) has dimension \(2N_f\),
  dimensional analysis then gives
  \[
    p=3+\frac{2N_f}{N_c-N_f}
      =\frac{3N_c-N_f}{N_c-N_f}.
  \]
  The anomalous axial symmetry fixes the same combination of \(\Lambda\)
  and \(M\).  Thus
  \[
    f_{\rm np}
      =C\left(
        \frac{\Lambda^{3N_c-N_f}}{\det M}
       \right)^{1/(N_c-N_f)}.
  \]
  Holomorphic decoupling fixes \(C=N_c-N_f\) in the standard
  normalization.

  \SupplementarySolution{20}{Holomorphic Decoupling of One Flavor}
  Split the meson into a light block \(\widehat M\) and the heavy diagonal
  entry \(X=M_{N_fN_f}\), setting mixed entries to zero at the stationary
  point.  With \(d=N_c-N_f\),
  \[
    f=d\left(
      \frac{\Lambda^{3N_c-N_f}}
           {X\det\widehat M}
      \right)^{1/d}+mX.
  \]
  The \(X\) equation eliminates it and produces
  \[
    f_{\rm light}
      =(d+1)\left(
       \frac{m\Lambda^{3N_c-N_f}}
            {\det\widehat M}
      \right)^{1/(d+1)}.
  \]
  This is exactly the ADS form for \(N_f-1\) flavors if
  \[
    \Lambda_{N_f-1}^{\,3N_c-(N_f-1)}
      =m\Lambda_{N_f}^{\,3N_c-N_f}.
  \]
  The matching is holomorphic and contains the complex phase of \(m\),
  not merely its magnitude.

  \SupplementarySolution{21}{Gaugino Condensation and Vacuum Multiplicity}
  Classically the gaugino phase rotation is a continuous \(U(1)_R\).
  An instanton has \(2N_c\) gaugino zero modes, so the anomaly preserves
  only \(\mathbb Z_{2N_c}\).  Strong dynamics generates
  \[
    \langle\lambda\lambda\rangle_k
      =C\Lambda^3e^{2\pi ik/N_c},
    \qquad k=0,\ldots,N_c-1.
  \]
  The bilinear has \(R=2\), so its expectation value breaks
  \(\mathbb Z_{2N_c}\) to the fermion-parity subgroup
  \(\mathbb Z_2\), leaving \(N_c\) inequivalent phases.  Each is a
  bosonic zero-energy vacuum.  Their net contribution to
  \(\operatorname{Tr}(-1)^F\) is therefore \(N_c\), agreeing with the
  Witten index and proving that supersymmetry is not dynamically broken
  in pure \(SU(N_c)\) super Yang--Mills theory.

  \SupplementarySolution{22}{Mass Deformation of the Quantum Constraint}
  Introduce
  \[
    f=A\left(\det M-B\bar B-\Lambda_H^{2N_c}\right)
      +mM^{N_c}{}_{N_c}.
  \]
  On the branch relevant after the last flavor is made heavy,
  \(B=\bar B=0\), the mixed mesons vanish, and
  \(\det M=M^{N_c}{}_{N_c}\det\widehat M\).  The constraint gives
  \[
    M^{N_c}{}_{N_c}
      =\frac{\Lambda_H^{2N_c}}{\det\widehat M}.
  \]
  Substitution leaves
  \[
    f_{\rm light}
      =\frac{m\Lambda_H^{2N_c}}{\det\widehat M}
      =\frac{\Lambda_L^{2N_c+1}}{\det\widehat M},
  \]
  where
  \(\Lambda_L^{2N_c+1}=m\Lambda_H^{2N_c}\).
  This is the ADS superpotential for \(N_f=N_c-1\), including its
  correct dimension and holomorphic threshold dependence.

  \SupplementarySolution{23}{Orthogonal-Group Holomorphic Dynamics}
  For \(SO(N_c)\), the adjoint index is \(N_c-2\), while each vector has
  index one.  Thus
  \[
    b_0=3(N_c-2)-N_f.
  \]
  If every vector chiral superfield has scalar \(R\) charge \(r\), the
  mixed gauge anomaly is
  \[
    (N_c-2)+N_f(r-1).
  \]
  It vanishes for
  \[
    r=1-\frac{N_c-2}{N_f}
      =\frac{N_f-N_c+2}{N_f}.
  \]
  The symmetric meson \(M_{ij}=Q_i\mathbin{\cdot}Q_j\) is the relevant
  flavor tensor.  Symmetry and dimension then uniquely give, for
  \(N_f<N_c-2\),
  \[
    f_{\rm np}
      =C\left(
       \frac{\Lambda^{3(N_c-2)-N_f}}{\det M}
      \right)^{1/(N_c-N_f-2)}.
  \]
  Its negative power of \(M\) drives the supersymmetric configuration
  to infinite field values, producing a runaway.

  \SupplementarySolution{24}{Broken Vacuum or Runaway?}
  For chiral fields with smooth positive Kahler metric \(g_{i\bar j}\),
  eliminating auxiliary fields gives
  \[
    V=g^{i\bar j}f_i f_{\bar j}^*
      +\frac12\sum_A D_A^2\geq0.
  \]
  At a finite stationary point, failure of all equations
  \(f_i=D_A=0\) means \(V>0\).  The supersymmetry algebra then implies
  that the vacuum is not annihilated by the supercharges, so
  supersymmetry is spontaneously broken and a goldstino exists.
  This conclusion presupposes a normalizable ground state.  If instead
  there is a path with fields tending to infinity and \(V\to0\), the
  infimum is zero but is not attained at finite distance.  The system
  has no vacuum state at all; it is a runaway, not a stable
  supersymmetry-breaking vacuum.  This distinction is also why an index
  may change when large-field asymptotics change.

  \SupplementarySolution{25}{O'Raifeartaigh Pseudomodulus and Goldstino}
  The three \(\mathcal F\)-flatness conditions are
  \[
    f_X=f_0+\frac h2\phi_1^2,\qquad
    f_{\phi_1}=hX\phi_1+m\phi_2,\qquad
    f_{\phi_2}=m\phi_1.
  \]
  The last equation forces \(\phi_1=0\); then the first cannot vanish
  when \(f_0\ne0\), while the second sets \(\phi_2=0\).  Therefore
  \[
    \phi_1=\phi_2=0,\qquad X\ {\rm arbitrary},\qquad
    V_{\rm tree}=\lvert f_0\rvert^2.
  \]
  The condition \(m^2>\lvert hf_0\rvert\) keeps both real scalar
  eigenvalues \(m^2\pm\lvert hf_0\rvert\) positive at the origin.
  The complex scalar \(X\) is a tree-level pseudomodulus.  Only
  \(\mathcal F_X=-f_0^*\) has a vacuum expectation value, so the
  normalized goldstino is simply the Majorana fermion \(\psi_X\), up to
  an overall phase.

  \SupplementarySolution{26}{Radiative Lifting of a Pseudomodulus}
  The one-loop potential for a constant \(X\) is
  \[
    V_1(X)=\frac{1}{64\pi^2}
      \operatorname{STr}\left[
       \mathcal M^4(X)
       \left(\ln\frac{\mathcal M^2(X)}{\mu^2}-\frac32\right)
      \right].
  \]
  Supersymmetry breaking splits the two bosonic eigenvalues associated
  with \(\phi_1,\phi_2\) from their fermionic partners by quantities of
  order \(hf_0\), while the \(X\)-dependent supersymmetric masses vary
  through \(hX\).  Expanding the supertrace near \(X=0\) gives
  \[
    m_X^2
      \sim\frac{\lvert h\rvert^4}{16\pi^2}
       \frac{\lvert f_0\rvert^2}{\lvert m\rvert^2}
       \,\mathcal G\!\left(
        \frac{\lvert hf_0\rvert}{\lvert m\rvert^2}
       \right),
  \]
  where the basic stable model has a positive dimensionless function
  \(\mathcal G\).  Thus the scalar pseudomodulus is lifted
  radiatively.  The fermion \(\psi_X\) remains exactly massless because
  the broken supercurrent Ward identity protects the goldstino, even
  though no supersymmetry pairs it with the now-massive scalar.

  \SupplementarySolution{27}{A Generic \(R\)-Symmetry Counting Test}
  At an \(R\)-symmetric point, all fields with nonzero \(R\) charge have
  zero expectation value.  A generic superpotential can then be written
  to leading order as
  \[
    f=\sum_{i=1}^{N_X}X_i\,g_i(Y_1,\ldots,Y_{N_Y})
      +\text{terms containing other charged fields},
  \]
  where \(R(X_i)=2\) and \(R(Y_a)=0\).  All auxiliary equations except
  those of the \(X_i\) vanish automatically at the symmetric locus, and
  the remaining conditions are
  \[
    g_i(Y)=0,\qquad i=1,\ldots,N_X.
  \]
  For generic functions, more independent equations than variables
  have no common solution.  Hence \(N_X>N_Y\) generically implies
  supersymmetry breaking at an \(R\)-symmetric vacuum.  When
  \(N_X\leq N_Y\), solutions generically exist; special coefficients,
  non-generic superpotentials, runaways, or vacua that spontaneously
  break \(R\) require a separate analysis.

  \SupplementarySolution{28}{The \(3\)-\(2\) Model Rank Obstruction}
  Denote the two \(SU(2)\) components of the \(SU(3)\) triplet \(Q\) by
  two flavors.  Strong \(SU(3)\) dynamics generates
  \[
    f_{\rm ADS}
      =\frac{\Lambda_3^7}
       {\det(\bar Q Q)}.
  \]
  The full superpotential is
  \(f=\lambda Q\bar D L+f_{\rm ADS}\).  The equation obtained by varying
  \(L\) is
  \[
    Q\bar D=0.
  \]
  Thus one column of the \(2\times2\) meson matrix
  \((\bar Q Q)\) vanishes, so its determinant is zero.  At precisely
  that locus the ADS term and its derivatives are singular; the
  remaining \(\mathcal F\)-term equations cannot vanish.  Moving toward
  large fields can reduce some \(\mathcal F\)-terms, but the
  \(SU(3)\) and \(SU(2)\) \(D\)-flatness conditions do not permit all
  dangerous directions simultaneously.  Their potential stabilizes the
  fields at finite values, leaving positive vacuum energy and dynamical
  supersymmetry breaking.

  \SupplementarySolution{29}{Quantum-Constraint Rank Condition}
  The confined coordinates are the antisymmetric mesons
  \(M_{ij}=Q_iQ_j\), constrained by
  \[
    \operatorname{Pf}M=\Lambda^{2N+2}.
  \]
  The singlet coupling gives
  \[
    \mathcal F_{S^{ij}}^*
      =-\lambda M_{ij}.
  \]
  Unbroken supersymmetry would therefore require every component of
  \(M\) to vanish.  That is impossible because its Pfaffian is fixed to
  a nonzero value by the exact quantum constraint.  Hence at least one
  singlet auxiliary field has a nonzero expectation value and the
  vacuum energy is positive.  A non-anomalous \(R\) symmetry assigns
  \(R(S)=2\) and \(R(M)=0\); it forbids additional generic terms that
  could remove the rank obstruction.  The goldstino is the normalized
  linear combination of singlet fermions weighted by their nonzero
  \(\mathcal F_S\) expectation values.

  \SupplementarySolution{30}{Current Normalization of the Goldstino Scale}
  Lorentz invariance and current conservation allow the leading
  vacuum-to-one-goldstino matrix element to be written
  \[
    \langle{\rm VAC}\lvert S^\mu(0)
      \rvert\mathbf p,\lambda\rangle
      =(2\pi)^{-3/2}
       \left[
        \frac{1+\gamma_5}{2}\gamma^\mu F_G
        +\frac{1-\gamma_5}{2}\gamma^\mu F_G^*
       \right]u(\mathbf p,\lambda)
  \]
  up to terms proportional to \(p^\mu\).  Inserting the two helicity
  states into the vacuum expectation value of the supercharge
  anticommutator and using
  \(\sum_\lambda u\bar u=-i\sl{p}/(2p^0)\) gives
  \[
    \rho_{\rm VAC}=\frac{\lvert F_G\rvert^2}{2}.
  \]
  The goldstino residues in matter fermions and gauginos are proportional
  to \(\sqrt2\langle\mathcal F_n\rangle/F_G\) and
  \(-\langle D_A\rangle/F_G\), respectively.  Normalizing the one-particle
  state therefore yields
  \[
    \lvert F_G\rvert^2
      =2\sum_n\lvert\langle\mathcal F_n\rangle\rvert^2
       +\sum_A\lvert\langle D_A\rangle\rvert^2,
  \]
  reproducing the auxiliary-field expression for the positive vacuum
  energy.
\end{enumerate}
